\begin{pastexam}[label={px:ncku-114-isec-ds-01}]{成大智慧資安學程 114 學年度 單選題第 1 題}
（單選；每題 5 分，答錯倒扣 2 分）

\textbf{一、單選題 (80 分，每題 5 分，答錯到扣2分)}

A is a two dimension array, the locations of A(3, 2) and A(2, 3) are 1115 and 1207, respectively.
Assume that each element occupies two addresses. Please give the location of A(5, 10).

a)~1841 \qquad b)~1851 \qquad c)~1861 \qquad d)~1871 \qquad e)~None of the above.
\end{pastexam}

\begin{pastsoln}
答案是 d) 1871。

題目沒說以列為主還是以行為主，先用兩個已知位址判斷（資結書第~4.1~節）。 設陣列有 $R$ 列、$C$ 行，
每個元素占 $2$ 個位址。

\textbf{以列為主（row-major）}：$A(2, 3)$ 在 $A(3, 2)$ 前面（第 2 列整列在第 3 列之前），位址應該比較小，
但 $1207 > 1115$，不合。 用公式算也一樣：
$\mathrm{loc}(2, 3) - \mathrm{loc}(3, 2) = 2\bigl[(2 - 3)C + (3 - 2)\bigr] = 2(1 - C) = 92$，得 $C = -45$，不可能。

\textbf{以行為主（column-major）}：同一行的元素連在一起，
\[
\mathrm{loc}(2, 3) - \mathrm{loc}(3, 2) = 2\bigl[(3 - 2)R + (2 - 3)\bigr] = 2(R - 1) = 1207 - 1115 = 92,
\]
所以 $R = 47$。 再從 $A(3, 2)$ 算到 $A(5, 10)$：行差 $8$、列差 $2$，
\[
\mathrm{loc}(5, 10) = 1115 + 2\bigl[(10 - 2) \times 47 + (5 - 3)\bigr] = 1115 + 2 \times 378 = 1115 + 756 = 1871 .
\]
下界（從 $0$ 還是 $1$ 開始）不影響，因為只用到兩個位置的差。 答案 d)。
\end{pastsoln}
